%============================================================
% CH2 · P2 · PROBLEM 1 — WHITE-BLACK BALANCED SUBTREES (CF 1676G)
%============================================================
\subsection{Problem 2.2.1: White-Black Balanced Subtrees}
\problemlabel{2.2.1}{ssec:ch2-white-black}

% sample 1 tree; vertex colours W / B, per-vertex style override g2p2-s<k>
\tikzset{
    g2p2-W/.style={gv},
    g2p2-B/.style={gv, fill=darkgray, text=white, draw=black},
    g2p2-s1/.style=g2p2-W, g2p2-s2/.style=g2p2-B, g2p2-s3/.style=g2p2-B, g2p2-s4/.style=g2p2-W,
    g2p2-s5/.style=g2p2-W, g2p2-s6/.style=g2p2-B, g2p2-s7/.style=g2p2-W,
    g2p2-S/.pic={
        \node[g2p2-s1] (s1) at (0,0) {1};
        \node[g2p2-s2] (s2) at (-1.2,-1) {2}; \node[g2p2-s3] (s3) at (1.2,-1) {3};
        \node[g2p2-s4] (s4) at (-1.2,-2) {4}; \node[g2p2-s5] (s5) at (0.6,-2) {5}; \node[g2p2-s6] (s6) at (1.8,-2) {6};
        \node[g2p2-s7] (s7) at (0.6,-3) {7};
        \foreach \a/\b in {1/2, 1/3, 2/4, 3/5, 3/6, 5/7} \draw[gtree] (s\a) -- (s\b);
    }
}

\problemcard{2.2.1}{White-Black Balanced Subtrees}{Codeforces 1676G}{https://codeforces.com/problemset/problem/1676/G}{\exeasy\ \cf{1300}\ (the purest subtree aggregation)}{Post-order aggregation of one number per vertex, $+1$ for white and $-1$ for black; a subtree is balanced iff its sum is 0; parents $a_i < i$ give the order for free}

\subsubsection*{1. Problem statement}

A tree with $n$ vertices is rooted at vertex 1; vertex $i \ge 2$ has parent $a_i$ with $a_i < i$. Each vertex is white (\texttt{W}) or black (\texttt{B}), given as a string $s$. The subtree of $u$ consists of $u$ and all its descendants. A subtree is \textbf{balanced} if it contains as many white as black vertices. Count the balanced subtrees (one per vertex, so the answer is the number of vertices whose subtree is balanced).

\textbf{Constraints:} $t \le 10^4$ test cases, $2 \le n \le 4000$, $\sum n \le 2\cdot10^5$, $1 \le a_i < i$.

\begin{table}[H]
\centering\small
\begin{tabular}{@{}llcl@{}}
\toprule
\textbf{$n$ / parents $a_2..a_n$} & \textbf{$s$} & \textbf{Output} & \textbf{Balanced subtrees} \\
\midrule
\texttt{7 / 1 1 2 3 3 5} & \texttt{WBBWWBW} & \texttt{2} & of 2 and 3 \\
\texttt{2 / 1} & \texttt{BW} & \texttt{1} & of 1 \\
\texttt{8 / 1 2 3 4 5 6 7} & \texttt{BWBWBWBW} & \texttt{4} & of 1, 3, 5, 7 \\
\bottomrule
\end{tabular}
\end{table}

\begin{diagrambox}
\begin{tikzpicture}[x=1cm, y=0.9cm]
    \begin{scope}
        \pic {g2p2-S};
        \node[font=\scriptsize\bfseries] at (0,0.8) {sample 1};
    \end{scope}
    \begin{scope}[xshift=4.4cm]
        \pic {g2p2-S};
        \node[draw=successgreen, line width=1.2pt, rounded corners=6pt, dashed, fit=(s2)(s4), inner sep=3pt] {};
        \node[draw=successgreen, line width=1.2pt, rounded corners=6pt, dashed, fit=(s3)(s5)(s6)(s7), inner sep=3pt] {};
        \node[font=\scriptsize\bfseries] at (0,0.8) {the two balanced subtrees};
    \end{scope}
    \node[note, anchor=west, text width=4.6cm, align=left] at (7.2,-1.4) {subtree of 2: $\{2, 4\}$: 1 W, 1 B\\subtree of 3: $\{3, 5, 6, 7\}$: 2 W, 2 B\\subtree of 1: 4 W, 3 B --- not balanced};
    \begin{scope}[yshift=-4.3cm, xshift=-1.2cm]
        \node[font=\scriptsize\bfseries, anchor=east] at (-0.5,0) {sample 3};
        \foreach \i in {1,...,8}{
            \ifodd\i \node[g2p2-B] (q\i) at (\i*1.25-1.25,0) {\i};
            \else \node[g2p2-W] (q\i) at (\i*1.25-1.25,0) {\i}; \fi
        }
        \foreach \i [evaluate=\i as \j using int(\i+1)] in {1,...,7} \draw[gtree] (q\i) -- (q\j);
        \foreach \i in {1,3,5,7} \node[gdist, below=1pt of q\i] {balanced};
    \end{scope}
\end{tikzpicture}
\figcap{Sample 1 (black disks are black vertices) and sample 3, a path rooted at 1 that alternates colours: the subtree of $i$ is the suffix $i..8$, balanced exactly when its length is even.}
\end{diagrambox}

\subsubsection*{2. The graph model}

\begin{modelbox}[Graph Model --- White-Black Balanced Subtrees]
\renewcommand{\arraystretch}{1.35}
\begin{tabularx}{\linewidth}{@{}>{\bfseries}l X@{}}
Vertices & the $n$ tree vertices, each with a weight $w(v) = +1$ (white) or $-1$ (black) \\
Edges & tree edges $a_i \to i$, given as a parent array; $a_i < i$, so every parent has a smaller label than its children \\
Weights & the vertex weights $\pm 1$ \\
Question & for how many $v$ is $\mathrm{bal}[v] = \sum_{x \in \mathrm{subtree}(v)} w(x)$ equal to 0? \\
Size & $\sum n \le 2\cdot10^5$: \bigO{n} per test; $n \le 4000$ even allows \bigO{n^2} per test --- but $t$ tests of $n = 4000$ cost $50 \cdot 4000^2 / 2 = 4\cdot10^8$ \\
\end{tabularx}
\end{modelbox}

\begin{diagrambox}
\begin{tikzpicture}[x=1cm, y=0.9cm]
    \begin{scope}
        \node[font=\scriptsize\bfseries, anchor=east] at (-0.6,0) {$i$};
        \foreach \i in {2,...,7} \node[cell, minimum width=8mm] (c\i) at (\i*0.9-1.8,0) {\i};
        \node[font=\scriptsize\bfseries, anchor=east] at (-0.6,-0.8) {$a_i$};
        \foreach \p [count=\i from 2] in {1,1,2,3,3,5} \node[cell, minimum width=8mm, fill=theorybg] at (\i*0.9-1.8,-0.8) {\p};
        \node[note, text width=5.2cm] at (2.2,-2.0) {every $a_i < i$: reading $i = n, n-1, \dots$ meets every child before its parent};
    \end{scope}
    \begin{scope}[xshift=7.6cm]
        \pic {g2p2-S};
        \foreach \v/\w/\d in {1/+1/right, 2/-1/left, 3/-1/right, 4/+1/left, 5/+1/left, 6/-1/right, 7/+1/right}
            \node[gdist, text=deepblue, \d=1pt of s\v] {$\w$};
        \node[note, anchor=west, text width=3.1cm, align=left] at (2.6,-1.5) {white $= +1$, black $= -1$:\\balanced $\iff$ sum $= 0$};
    \end{scope}
\end{tikzpicture}
\figcap{The model: the parent array is already a valid order, and one signed number per vertex replaces the two colour counts.}
\end{diagrambox}

\subsubsection*{3. Key observations}

\begin{itemize}[leftmargin=*]
    \item \textbf{One number instead of two.} ``$\#W = \#B$'' is ``$\#W - \#B = 0$''. With $w = \pm 1$ the subtree sum $\mathrm{bal}[v]$ is all we need, and it aggregates: $\mathrm{bal}[v] = w(v) + \sum_c \mathrm{bal}[c]$.
    \item \textbf{The labels are an order.} Since $a_i < i$, all children of $v$ have labels $> v$. Looping $v = n, n-1, \dots, 1$ processes every child before its parent: at the moment we reach $v$, $\mathrm{bal}[v]$ is final (the pattern's invariant) and can be tested and handed up.
    \item \textbf{No graph is built.} No adjacency lists, no search, no recursion: two arrays.
    \item \textbf{Odd subtrees are never balanced.} A balanced subtree has even size --- a quick sanity check on traces (in sample 1 the balanced ones have sizes 2 and 4).
\end{itemize}

\subsubsection*{4. Approaches}

\begin{bruteforcebox}
\textbf{Idea: list every subtree and count.} For each $u$, collect its subtree by walking only from parents to children, and compare the two counts.

\begin{lstlisting}[style=cpp]
#include <bits/stdc++.h>
using namespace std;

// Brute force: for every vertex u, list its subtree explicitly (a search that only walks
// from a vertex to its children) and count the two colours. O(n^2) per test.
int main() {
    int t;
    cin >> t;
    while (t--) {
        int n;
        cin >> n;
        vector<vector<int>> children(n + 1);
        for (int i = 2; i <= n; i++) {
            int p;
            cin >> p;
            children[p].push_back(i);
        }
        string s;
        cin >> s;
        int balanced = 0;
        for (int u = 1; u <= n; u++) {
            int white = 0, black = 0;
            vector<int> todo(1, u);                // the subtree of u = everything below u
            while (!todo.empty()) {
                int x = todo.back(); todo.pop_back();
                if (s[x - 1] == 'W') white++; else black++;
                for (int c : children[x]) todo.push_back(c);
            }
            if (white == black) balanced++;
        }
        cout << balanced << "\n";
    }
}
\end{lstlisting}

\begin{complexitybox}
\textbf{Time:} $\sum_u \mathrm{sz}[u]$ = $\sum_v (\mathrm{dep}[v] + 1)$, which is $\Theta(n^2)$ on a path. Fifty paths of $4000$ vertices ($\sum n = 2\cdot10^5$) took about 3\,s here, against $0.26$\,s for Approach~2 including process start-up. \textbf{Space:} \bigO{n}. \textbf{Verdict:} borderline at these limits, hopeless if $n$ were $2\cdot10^5$ --- and the reference for testing.
\end{complexitybox}
\end{bruteforcebox}

\begin{algorithmbox}[Approach 2: one backward pass over the labels]
\textbf{Idea.} $\mathrm{bal}[v] = \pm1$ initially. For $v = n$ down to $1$: $\mathrm{bal}[v]$ is now final; count it if it is 0; add it to $\mathrm{bal}[a_v]$. This is template (3) of the pattern with $\pm1$ in place of $1$.

\begin{lstlisting}[style=cpp]
#include <bits/stdc++.h>
using namespace std;

// White-Black Balanced Subtrees (CF 1676G). Parents satisfy a_i < i, so the labels n, n-1, ..., 2
// already list every child before its parent: bal[v] = (#white - #black) in the subtree of v.
int main() {
    ios::sync_with_stdio(false);
    cin.tie(nullptr);
    int t;
    cin >> t;
    while (t--) {
        int n;
        cin >> n;
        vector<int> par(n + 1, 0);
        for (int i = 2; i <= n; i++) cin >> par[i];
        string s;
        cin >> s;
        vector<int> bal(n + 1, 0);
        for (int v = 1; v <= n; v++) bal[v] = (s[v - 1] == 'W') ? 1 : -1;
        int balanced = 0;
        for (int v = n; v >= 1; v--) {             // children (larger labels) are finished first
            if (bal[v] == 0) balanced++;           // the subtree of v is complete: test it
            if (v > 1) bal[par[v]] += bal[v];      // hand the finished value to the parent
        }
        cout << balanced << "\n";
    }
}
\end{lstlisting}

\begin{complexitybox}
\textbf{Time:} \bigO{n} per test, \bigO{\sum n} in total. \textbf{Space:} \bigO{n}. \textbf{Verdict:} reproduces all samples; no recursion, so depth is irrelevant.
\end{complexitybox}
\end{algorithmbox}

\begin{invariantbox}[Why the test happens at the right moment]
\textbf{Invariant.} When the loop reaches $v$, $\mathrm{bal}[v] = \sum_{x \in \mathrm{subtree}(v)} w(x)$.

Every descendant $x$ of $v$ has a label $> v$ (labels grow along every root path because $a_i < i$), so it was processed earlier, and when it was processed it added its own final value to its parent. By induction from the largest label down, each child $c$ of $v$ added exactly $\mathrm{bal}[c] = $ the sum over its subtree, and nothing else was ever added to $\mathrm{bal}[v]$ (only children write into it). The children's subtrees are disjoint and together with $v$ form the subtree of $v$. Hence the value is final and complete when $v$ is reached, and testing it there counts each balanced subtree exactly once. \hfill$\blacksquare$
\end{invariantbox}

\begin{alternativebox}[Approach 3: a real post-order DFS]
Build children lists and let a recursive DFS return the subtree balance. The depth is at most $n \le 4000$ frames --- far below the $\approx 65\,000$ that a 2\,MB stack survives (Lesson~\ref{ssec:ch2-call-stack}), so recursion is safe \emph{here}. This version does not depend on $a_i < i$ and works for any parent array.

\begin{lstlisting}[style=cpp]
#include <bits/stdc++.h>
using namespace std;

// Alternative: a real post-order DFS on the children lists. It returns (#white - #black)
// of the subtree and counts balanced subtrees on the way out. Depth <= n <= 4000: safe.
vector<vector<int>> children;
string s;
int balanced;

int dfs(int u) {
    int bal = (s[u - 1] == 'W') ? 1 : -1;
    for (int c : children[u]) bal += dfs(c);  // every child subtree is finished here
    if (bal == 0) balanced++;
    return bal;
}

int main() {
    ios::sync_with_stdio(false);
    cin.tie(nullptr);
    int t;
    cin >> t;
    while (t--) {
        int n;
        cin >> n;
        children.assign(n + 1, vector<int>());
        for (int i = 2; i <= n; i++) {
            int p;
            cin >> p;
            children[p].push_back(i);
        }
        cin >> s;
        balanced = 0;
        dfs(1);
        cout << balanced << "\n";
    }
}
\end{lstlisting}
\begin{complexitybox}
\textbf{Time:} \bigO{n} per test. \textbf{Space:} \bigO{n} plus a stack of height $\le n$. \textbf{Verdict:} reproduces the samples; agrees with Approaches 1 and 2 on all random tests.
\end{complexitybox}
\end{alternativebox}

\begin{tricksbox}[Approach 4: Implementation Refinements]
\begin{itemize}[leftmargin=*]
    \item \textbf{Reset per test.} With $t = 10^4$ tests, allocate arrays of size $n + 1$ per test (as above), never of size $2\cdot10^5$ per test --- that alone would be $2\cdot10^9$ writes.
    \item \textbf{Two counters work too.} Keeping \texttt{white[v]} and \texttt{black[v]} is equally correct; the signed sum halves the memory and is the form that generalises (``sum of $\pm1$ weights $= 0$'').
    \item \textbf{The root is a subtree too.} The loop includes $v = 1$; sample 2 (\texttt{BW}) has its only balanced subtree at the root.
\end{itemize}
\end{tricksbox}

\subsubsection*{5. Trace}

\begin{tracebox}[{Sample 1: $a = (1,1,2,3,3,5)$, $s = \texttt{WBBWWBW}$, initial $\mathrm{bal} = (+1,-1,-1,+1,+1,-1,+1)$}]
\begin{center}
\begin{tikzpicture}[x=0.85cm, y=0.8cm]
    \tikzset{gv/.append style={minimum size=6.5mm, font=\scriptsize\bfseries}}
    \begin{scope}
        \tikzset{g2p2-s7/.style={g2p2-W, draw=tabulationblue, line width=1.3pt}, g2p2-s6/.style={g2p2-B, draw=tabulationblue, line width=1.3pt},
                 g2p2-s5/.style={g2p2-W, draw=accentorange, line width=1.6pt}}
        \pic {g2p2-S};
        \foreach \v/\b/\d in {1/+1/right, 2/-1/left, 3/-2/right, 4/+1/left, 5/+2/left, 6/-1/right, 7/+1/right} \node[gdist, text=deepblue, \d=0pt of s\v] {$\b$};
        \node[font=\scriptsize\bfseries] at (0,0.9) {1. $v = 7, 6$ done; at $v = 5$};
    \end{scope}
    \begin{scope}[xshift=5.4cm]
        \tikzset{g2p2-s3/.style={g2p2-B, draw=successgreen, line width=1.8pt}, g2p2-s2/.style={g2p2-B, draw=accentorange, line width=1.6pt}}
        \pic {g2p2-S};
        \foreach \v/\b/\d in {1/+1/right, 2/0/left, 3/0/right, 4/+1/left, 5/+2/left, 6/-1/right, 7/+1/right} \node[gdist, text=deepblue, \d=0pt of s\v] {$\b$};
        \node[font=\scriptsize\bfseries] at (0,0.9) {2. after $v = 5, 4$: 3 and 2 at 0};
    \end{scope}
    \begin{scope}[xshift=10.8cm]
        \tikzset{g2p2-s3/.style={g2p2-B, draw=successgreen, line width=1.8pt}, g2p2-s2/.style={g2p2-B, draw=successgreen, line width=1.8pt},
                 g2p2-s1/.style={g2p2-W, draw=accentorange, line width=1.6pt}}
        \pic {g2p2-S};
        \foreach \v/\b/\d in {1/+1/right, 2/0/left, 3/0/right, 4/+1/left, 5/+2/left, 6/-1/right, 7/+1/right} \node[gdist, text=deepblue, \d=0pt of s\v] {$\b$};
        \node[font=\scriptsize\bfseries] at (0,0.9) {3. end: count $= 2$};
    \end{scope}
\end{tikzpicture}
\end{center}
\begin{center}\small
\begin{tabular}{@{}cccl@{}}
\toprule
$v$ & $\mathrm{bal}[v]$ (final) & balanced? & hand up \\
\midrule
7 & $+1$ & no & $\mathrm{bal}[5]$: $+1 \to +2$ \\
6 & $-1$ & no & $\mathrm{bal}[3]$: $-1 \to -2$ \\
5 & $+2$ & no & $\mathrm{bal}[3]$: $-2 \to 0$ \\
4 & $+1$ & no & $\mathrm{bal}[2]$: $-1 \to 0$ \\
3 & $0$ & \textbf{yes} (count 1) & $\mathrm{bal}[1]$: $+1 \to +1$ \\
2 & $0$ & \textbf{yes} (count 2) & $\mathrm{bal}[1]$: $+1 \to +1$ \\
1 & $+1$ & no & --- \\
\bottomrule
\end{tabular}
\end{center}
Printed by a copy of Approach~2 with one trace line per step. Mistake 1 below was also run: the forward loop prints \texttt{1}. Orange ring: the vertex being processed; green ring: a balanced subtree; blue rings: already handed up.
\end{tracebox}

\subsubsection*{6. Test suite}

\begin{testbox}
\centering\small
\begin{tabular}{@{}p{5.2cm}cp{6.4cm}@{}}
\toprule
\textbf{Input ($n$ / parents / $s$)} & \textbf{Output} & \textbf{Why} \\
\midrule
\texttt{2 / 1 / WW} & \texttt{0} & no subtree can balance \\
\texttt{2 / 1 / BW} & \texttt{1} & only the root (sample 2) \\
\texttt{5 / 1 1 1 1 / WBBBB} & \texttt{0} & star: leaves have size 1, root has 1 W vs 4 B \\
\texttt{4 / 1 1 1 / BWBW} & \texttt{1} & star: only the root \\
\texttt{6 / 1 2 3 4 5 / WBWBWB} & \texttt{3} & path: the suffixes of even length \\
50 paths, $n = 4000$, \texttt{WBWB\dots} & \texttt{2000} each & largest input; deepest tree \\
\bottomrule
\end{tabular}
\end{testbox}
All outputs from Approach~2; Approach~1 printed the same. In addition, Approaches 1--3 agreed on $3000$ random inputs (1--3 test cases each, $n \le 12$, random / path / star / caterpillar shapes).

\subsubsection*{7. Common mistakes}

\begin{warningbox}[Mistake 1 --- the forward loop]
Looping $v = 2, \dots, n$ and adding $\mathrm{bal}[v]$ to the parent hands up values \emph{before} their own children have arrived. On sample 1, vertex 5 would hand up $+1$ instead of $+2$, $\mathrm{bal}[3]$ would end at $-1$, and the answer would be 1 instead of 2.
\end{warningbox}

\begin{warningbox}[Mistake 2 --- testing before the subtree is complete]
Counting zeros \emph{while} adding (``if \texttt{bal[par[v]]} became 0, count it'') counts a parent that passes through 0 only temporarily. On \texttt{3 / 1 1 / BBW} the root goes $-1 \to 0$ (vertex 3 arrives) $\to -1$ (vertex 2 arrives): that version prints \texttt{1}, the correct answer is \texttt{0}. Test each $v$ once, at the moment it is reached.
\end{warningbox}

\begin{warningbox}[Mistake 3 --- off-by-one in the string]
Vertices are 1-based, the string is 0-based: the colour of $v$ is \texttt{s[v-1]}.
\end{warningbox}

\subsubsection*{8. Takeaways}

\begin{itemize}[leftmargin=*]
    \item Turn ``two counts equal'' into ``a signed sum is 0'': one aggregated number.
    \item A parent array with $a_i < i$ is an order: one backward loop is a post-order traversal.
    \item Test a subtree's value exactly once, when the vertex is left (reached in the backward loop).
\end{itemize}

\begin{variantsbox}
\variant{Count per label}{\varin{https://leetcode.com/problems/number-of-nodes-in-the-sub-tree-with-the-same-label/}}
  {for each vertex, how many vertices in its subtree have its own label (26 labels).}
  {the same tree, a counter per label instead of one signed number.}
  {a global counter, read at entry and at exit of $v$ (Exercise E2.1).}
\variant{Three colours, all equal}{\varnone}
  {colours R, G, B; a subtree is balanced if the three counts are equal.}
  {two signed sums: $\#R - \#G$ and $\#G - \#B$.}
  {aggregate the pair; balanced iff both are 0. Matches an explicit-count brute force on 5000 random trees.}
\variant{Arbitrary edge list}{\varnone}
  {the tree is given as $n - 1$ undirected edges, root 1.}
  {no label order any more.}
  {template (2): BFS order, backward loop. Matches the label loop on 5000 randomly relabelled trees.}
\variant{Even subtrees to cut}{\varin{https://codeforces.com/problemset/problem/982/C}}
  {cut as many edges as possible so that every piece has even size.}
  {aggregate sizes; an edge above $v$ can be cut iff $\mathrm{sz}[v]$ is even.}
  {the same backward loop with a parity test (Exercise E2.2).}
\end{variantsbox}
